\documentclass[11pt]{article}
\usepackage[T1]{fontenc}
\usepackage[margin=1in]{geometry}
\usepackage{amsmath,amssymb,amsthm,microtype}
\usepackage[colorlinks=true,allcolors=blue]{hyperref}
\newtheorem{theorem}{Theorem}
\newtheorem{lemma}[theorem]{Lemma}
\title{Query complexity of Boolean MCSP with polynomially separated size thresholds}
\author{Research manuscript}\date{22 November 2026}
\begin{document}\maketitle
\begin{abstract}
A reader can query entries of an $N$-bit truth table. The promise is that its Boolean circuit size is at most $s$ or greater than $t$. For $s=\Theta(N^a)$ and $t=\Theta(N^b)$, with fixed $0<a<b<1$, we show that the bounded-error cost is $\Theta(Ns/t)$ classical queries and $\Theta(s\sqrt{N/t})$ quantum queries. The hard tables have one mark per bucket. Low-degree polynomials over a binary extension field place the marks of easy tables in positions that look independent to a reader who has seen only a small number of them. A Boolean circuit for evaluating any one of these field polynomials fits the smaller size threshold. We also show that the threshold degree, globally bounded rational approximate degree and bounded-error postselected quantum query complexity of this promise are all of order $s$.
\end{abstract}

\section{Estimating circuit size with a promise}
Fix an inner Boolean basis containing AND, OR and NOT, with binary fan-in, free constants and unrestricted fanout. Let $C(f)$ count gates. Put $N=2^n$ and write $A_u$ for tables representing functions on $n$ variables with $C(f)\le u$. A tester must accept every table in $A_s$ and reject every table outside $A_t$. We count truth-table queries, allow unlimited computation between them, and require error at most $1/3$ on each promised input. Logs are base two.
\begin{theorem}
If $s=\Theta(N^a)$ and $t=\Theta(N^b)$ for fixed $0<a<b<1$, the classical randomized and quantum query complexities satisfy
\[
                 R=\Theta(Ns/t),\qquad Q=\Theta(s\sqrt{N/t}).
\]
The upper bounds allow perfect completeness. The implicit constants may depend on the exponents and the constants in the size thresholds.
\end{theorem}
Here the gap is a promise on the minimum number of gates. Hamming distances enter only in the proof of the upper bound.
\section{Upper bounds by testing distance}
For some constant $c_b>0$, every Boolean table of weight at most $c_b t$ has a circuit with fewer than $t/2$ gates for large $n$. Indeed Lupanov's sparse synthesis bound \cite{kors} at weight $\lceil c_bt\rceil$ costs $(1+o(1))E/\log E$, where $E=\log\binom{2N}{\lceil c_bt\rceil}=(c_b(1-b)+o(1))tn$ and $\log E=(b+o(1))n$. An extra address bit pads a lower weight to this weight and can be fixed after synthesis. The theorem's weight condition $t/\log n\to\infty$ holds. XOR with a circuit of size at most $s=o(t)$ costs only a constant number of further gates. Thus any no-instance is at Hamming distance at least $c_b t$ from $A_s$.

We give an algorithm that accepts every yes input with certainty. Write $B=\log |A_s|=O(s\log s)$ by ordinary rooted circuit counting. Put $\xi=c_bt/(2N)$, $v=\xi$ classically and $v=\sqrt\xi$ quantumly. Maintain all size-$s$ candidates consistent with the verified answers. Choose a coordinate with minority fraction at least $\gamma=c v\log(2/v)$ if one exists, using a small absolute $c$. Query it and retain exactly consistent candidates. Such a query reduces the candidate count by a factor at most $1-\gamma$, so at most $O(B/\gamma)$ occur.

Otherwise let $m$ be the coordinatewise majority of the survivors. This table is close to one computable with few gates. Draw five independent uniform survivors and take their majority on each coordinate. It has circuit size at most $5s+O(1)$. Its expected disagreement with $m$ is $O(\gamma^3)N=o(\xi N)$, so one choice achieves this distance. On any no input, the actual table and $m$ must differ in at least $\xi N$ entries. Otherwise the input would lie within $c_b t$ entries of the five-circuit majority and could be synthesized with fewer than $t$ gates, by the sparse correction bound and $5s=o(t)$.

Search for an actual disagreement with $m$, repeating independent random samples or bounded-cost quantum search attempts until a verified witness appears. On a no input the expected cost is $O(1/v)$ per witnessed mismatch. Each successful mismatch retains at most a $\gamma$ fraction of survivors. It therefore takes at most $1+B/\log(1/\gamma)$ successes to exhaust the candidate set. Reject if the set becomes empty. Cap the computation at a large constant multiple of the no-input expected query bound and accept at timeout. No true size-$s$ circuit is ever discarded on a yes input. On a no input Markov's inequality bounds the timeout probability by $1/3$. The query bound is
\[
 O\!\left(\frac1v+\frac{B}{v\log(2/v)}\right).
\]
Here $\log(2/v)=\Theta(n)$, so substitution gives both upper bounds. The constant-five majority is where we use the growing ratio between the two thresholds.

\section{Encoding many marks with a short circuit}
Let $k$ be a power of two between $C_b t$ and $2C_b t$, for a sufficiently large constant $C_b$. Divide the table into $k$ buckets of size $H=N/k$. The table has a single mark at position $Y_i\in\{0,1\}^h$ in bucket $i$, where $h=\log H$. Put $\ell=\max(\log k,h)$ and identify binary $\ell$-tuples with a basis of $F=\mathbb F_{2^\ell}$. Assign each bucket its distinct field element by zero-padding its binary address. Choose a uniform polynomial of degree less than $D=\lfloor c_a s\rfloor$ in $F[X]$ and take the first $h$ bits of its evaluation for the location. Vandermonde interpolation gives jointly uniform labels in every $D$ buckets.

\begin{lemma}\label{eval}
Each fixed degree-below-$D$ polynomial on $F$ has a Boolean evaluation circuit of size
\[
                  O\bigl(D\ell/\log D+\sqrt D\,\ell^{O(1)}\bigr)
\]
for $D\to\infty$.
\end{lemma}
\begin{proof}
Take $m=\lceil\sqrt D\rceil$. Form $1,x,\ldots,x^{m-1},x^m$ and the $m$ giant powers $(x^m)^j$ using $O(m)$ field multiplications of polynomial cost in $\ell$. Pad the coefficients to an $m$ by $m$ array $c_{ji}$. The $j$th block result is $S_j=\sum_{i=0}^{m-1} c_{ji}x^i$.

Multiplication in any binary extension of degree $\ell$ has bilinear rank $r=O(\ell)$ \cite{rank}. In particular $uv=\sum_{e=1}^r w_e L_e(u)R_e(v)$ for fixed $w_e\in F$ and binary linear functionals $L_e,R_e$. For a fixed $e$, the $m$ bits $R_e(x^i)$ form one vector. Map it to the $m$ sums $\sum_i L_e(c_{ji})R_e(x^i)$ by one hardwired binary linear map. Every $m$ by $m$ binary map costs $O(m^2/\log m)$ XOR gates: group inputs into blocks of $\lfloor(\log m)/2\rfloor$ bits, generate all sums in each block, and combine one partial sum per block in each output. At all $r$ ranks together this is $O(D\ell/\log D)$. Forming the input linear forms and reconstructing the $S_j$ costs $O(m\ell^2)$. Finally combine $S_j$ with the giant powers and add the field products. Implementing each XOR by a constant number of basis gates leaves the bound unchanged.
\end{proof}
Since $\ell\le n$ and $\log D\sim an$, choosing $c_a>0$ sufficiently small leaves room to compare the computed $h$-bit label with the suffix of the truth-table address. Thus every code-generated table belongs to $A_s$. The independent comparison law chooses all $Y_i$ independently and uniformly. Its tables have point probability $2^{-kh}$ and $kh\ge(C_b(1-b)+o(1))tn$. Rooted counting gives $\log|A_t|\le(t+1)\log(t+n+2)+O(t+\log n)=(b+o(1))tn$. Choose $C_b>b/(1-b)$ with margin. A table from this law lies outside $A_t$ with probability tending to one.

\section{What a query can learn about the marks}
\paragraph{Classical queries.}
Fix a deterministic query tree of depth $u\le cDH$, for a sufficiently small absolute $c$. Repeated bit queries can be answered locally. The following argument compares its distributions of transcripts under independent and $D$-wise independent labels. For each leaf expose the full labels of every bucket with a positive answer and every bucket that received at least $H/2$ distinct queries. The latter number is at most $2u/H$. The resulting refined leaves partition the label space. Call a leaf good if at most $D/4$ buckets with a positive answer were encountered.

Under the independent law the expected number of first hits before a bucket reaches $H/2$ queries is at most $2u/H$. Conditional on its observed misses, its next hit probability is at most $2/H$. Any remaining positive buckets are among the at most $2u/H$ heavy buckets. Thus the bad probability is at most $16c$ by Markov. At a good refined leaf at most $D/4+2u/H<D/2$ labels have been exposed. The two laws give these exposed values the same probability; under their conditioning the other labels are still $(D/2)$-wise uniform, rounding down. On the unexposed buckets the leaf asks that $Y_i$ avoid some set of size $m_i<H/2$, with $\sum_i m_i\le u$. Conditional on independent labels the avoidance probability is at least $\exp(-2u/H)$. The first $j=\lfloor D/2\rfloor$ terms of inclusion--exclusion are common to both laws. Bonferroni bounds their difference by at most twice the $j$th elementary symmetric sum in the numbers $m_i/H$, hence by $2(eu/(jH))^j$. For small enough $c$ this is exponentially small in $D$ even after division by $\exp(-2u/H)$. Each good refined leaf thus has nearly the same relative probability under the two laws. Summing over good leaves controls the mass of the bad leaves in both laws by $16c+o(1)$. A randomized mixture of trees has the same indistinguishability bound. This proves $R=\Omega(DH)$.

\paragraph{Quantum queries.}
For completeness, view the quantum state as a vector-valued function of $Y$. Under the uniform law, decompose this function orthogonally by its sets of mean-zero bucket coordinates; the level is the size of such a set. For levels up to $j=\lfloor D/2\rfloor$, inner products agree under both laws, since each term in an inner product depends on at most $2j$ labels. A query increases the level by at most one. On a fixed address $(i,z)$ it acts by $I+\mathbf1[Y_i=z](X-I)$, where $X$ is the answer-bit flip. Its upward transition adds the mean-zero part of $\mathbf1[Y_i=z]$, of norm at most $H^{-1/2}$. Distinct supports constant in $i$ remain orthogonal after adding $i$. Orthogonal address blocks handle superposed queries, giving upward operator norm at most $u_0=2/\sqrt H$.

Project back to levels below $j$ after each query in an auxiliary calculation. This calculation agrees under the two laws. If $a_v(q)$ denotes its norm on levels at least $v$ after $q$ queries, then $a_v(q+1)\le a_v(q)+u_0a_{v-1}(q)$ and $a_0(q)\le1$. The norm discarded over $q$ queries is at most $\binom qj u_0^j$ by the binomial recurrence and telescoping. The true final states therefore have acceptance probabilities differing by at most $4\binom qj u_0^j$. At $q\le c' D\sqrt H$ this is $o(1)$ for a small constant $c'$. The promised labels require a constant gap, proving $Q=\Omega(D\sqrt H)$.

\section{Verifying an externally specified hard table}

Fix a table \(x\) of circuit size greater than \(t\). Its fractional certificate complexity \(\mathrm{FC}(x)\) is the minimum of \(\sum_i a_i\) over nonnegative weights on the \(N\) positions such that \(\sum_{i:y_i\ne x_i}a_i\ge1\) for every table \(y\) of size at most \(s\). Its block sensitivity on this promise is the maximum number of pairwise disjoint sets of positions such that flipping any one set in \(x\) yields a size-\(s\) table.

\begin{theorem}
Every hard table has block sensitivity exactly two if \(t\ge3s+O(1)\). Uniformly over all hard tables,
\[
   2\le \mathrm{FC}(x)\le 2+O\!\left(\sqrt{(s+1)\log N/t}\right)
   \qquad\text{if }t/((s+1)\log N)\to\infty .
\]
\end{theorem}
\begin{proof}
The two constant tables give disjoint blocks whose union is all positions. This proves both lower bounds. If three disjoint flips produce easy \(y_1,y_2,y_3\), their coordinatewise majority is \(x\), giving it a circuit of size at most \(3s+O(1)\).

Normalize fractional certificate weights to a distribution and apply finite minimax to the disagreement indicator. This gives
\[
 \frac1{\mathrm{FC}(x)}
 =\max_{\mu\ \mathrm{on\ positions}}\min_{y\ \mathrm{easy}}
         \Pr_{i\sim\mu}[y_i\ne x_i]
 =\min_{\nu\ \mathrm{on\ easy\ tables}}\max_i
         \Pr_{y\sim\nu}[y_i\ne x_i]=:p_* .
\]
The constants ensure \(p_*\le1/2\). Choose a minimizing \(\nu\) and sample an odd number \(r\) of independent easy tables from it. At each position the probability that their majority disagrees with \(x\) is at most
\([4p_*(1-p_*)]^{r/2}\) by the binomial Chernoff bound. Their majority costs at most \(r(s+C)\) gates for an absolute constant \(C\). For example, full adders reduce triples of equal-significance bits to a sum and a carry until only two binary words remain. One final addition and comparison complete the majority in \(O(r)\) gates. If \(r(s+C)\le t\), every sampled majority must err somewhere. A union bound therefore yields
\[
 p_*\ge\tfrac12\left(1-\sqrt{1-N^{-2/r}}\right).
\]
Taking the largest admissible odd \(r\) proves the assertion.
\end{proof}

For the polynomially separated budgets of this paper, an optimal distribution \(\mu\) for a \emph{known} hard table \(x\) catches each easy competitor with probability at least \(1/2-o(1)\) in one comparison against the specified bit of \(x\). Repetition gives constant-query verification of equality to this \(x\) under the promise that the alternative is easy. The verifier, including its sampling distribution, depends on the whole reference \(x\). The decision bounds proved above use one algorithm for all promised input tables.

\begin{theorem}
For \(s=\Theta(N^a), t=\Theta(N^b)\), \(0<a<b<1\) fixed, the nonnegative quantum adversary bound is \(\Theta(\sqrt{N/t})\). The classical minimax adversary bound and the maximum easy-side block sensitivity and fractional certificate complexity are all \(\Theta(N/t)\).
\end{theorem}
\begin{proof}
A smooth version of the preceding minimax argument gives each hard \(x\) a distribution \(\mu_x\) with
\[
 \max_i\mu_x(i)\le1/k,\qquad
 \min_{y\ {\rm easy}}\Pr_{i\sim\mu_x}[x_i\ne y_i]\ge1/2-o(1),
 \qquad k=\lfloor c_b t\rfloor .
\]
Here \(c_b\) is the small constant in the sparse synthesis estimate above. In the minimax game restrict the position player to distributions with the stated point-mass cap. Against a fixed mixture of easy tables its best payoff is the average of the \(k\) largest pointwise error probabilities. Call the minimum of this average over the mixtures \(p\). Outside the \(k\) worst positions every error probability of a minimizing mixture is at most \(p\). Take the largest odd \(r\) with \(r(s+C)\le t/3\). If \(p<1/2-O(\sqrt{n/r})\) for a sufficiently large absolute implied constant, some majority is exact everywhere outside those \(k\) positions by the same union bound. Correcting the remaining discrepancies by a circuit for a weight-at-most-\(k\) table costs less than \(t/2\), and an XOR would then compute \(x\) within \(t\) gates. This contradiction proves the claim. The smooth minimax argument is the hard-core distribution argument specialized to truth tables admitting cheap sparse corrections.

In the \href{https://theoryofcomputing.org/articles/v002a001/}{nonnegative adversary minimax formula}, assign the uniform position law to every easy table \(y\), and \(\mu_x\) to hard \(x\). For every opposite pair,
\[
 \sum_{i:x_i\ne y_i}\sqrt{\mu_x(i)/N}
 \ \ge\sqrt{k/N}\sum_{i:x_i\ne y_i}\mu_x(i)=\Omega(\sqrt{t/N}).
\]
The reciprocal bounds the nonnegative adversary from above. In the \href{https://drops.dagstuhl.de/entities/document/10.4230/LIPIcs.STACS.2018.8}{classical minimax formula}, the summand is the minimum of the two position masses. The uniform law on both inputs already gives denominator at least \(\Omega(t/N)\) since every easy--hard distance is \(\Omega(t)\). This separation of the two promised sets also bounds fractional certificate complexity at every easy table by \(O(N/t)\), using uniform weights.

For the lower bounds, randomly permute the positions and cut out disjoint blocks of \(m=\lceil C_b t\rceil\) positions for a sufficiently large constant \(C_b\). Each block is a uniform \(m\)-subset, so its indicator is easy even at budget \(t\) with probability at most \(|A_t|/\binom Nm=\exp(-\Omega(t n))\), by the census and \(m\log(N/m)\sim C_b(1-b)tn\) in base-two logs. A union bound fixes \(\lfloor N/m\rfloor\) disjoint blocks all having hard indicators. Around the easy zero table these give that many disjoint sensitive blocks, and hence the same lower bound for its fractional certificates. For arbitrary input-specific position laws, one of the blocks has mass \(O(t/N)\) under the law for zero. Paired with that block's hard indicator, the classical minimax denominator is at most that mass and the quantum denominator at most its square root by Cauchy--Schwarz. These give the matching adversary lower bounds.
\end{proof}

\paragraph{Integral certificates.}
In the same polynomially separated regime every hard input has minimum certificate size \((1+o(1))s\log_2s\). Here a certificate must exclude every promised input of the other type. Indeed, sample \(\lceil(1+\delta)s\log_2s\rceil\) positions independently from the above comparison law for the specified hard input. A fixed easy competitor survives with probability at most \((1/2+o(1))^{(1+\delta)s\log_2s}\). The census and a union bound leave no survivor for some sample, and repeated positions can be deleted. Conversely every partial table on at most \((1-\delta)s\log_2s\) positions extends to an easy table. Apply \href{https://keldysh.ru/papers/2024/mvk/mvk2024_152.pdf}{Chashkin's Theorem~6.1} on arbitrary prescribed domains, allowing up to half of the entries to have the minority label. For a domain of size \(d=\lfloor(1-\delta)s\log_2s\rfloor\) its extension cost is at most \((1+o(1))d/\log_2d+O(n)<s\); smaller domains can first be padded to \(d\).

Every easy input has minimum certificate size \(N-(1+o(1))t\log_2t\). If a certificate leaves \(m\) positions unqueried, all \(2^m\) completions must have size at most \(t\), so \(m\le\log_2|A_t|\le(1+o(1))t\log_2t\). Conversely leave unqueried the first \(\lfloor(1-\delta)t\log_2t\rfloor\) addresses. Modifications on that interval are arbitrary binary vectors of this length and zero outside. Dyadic prefix decomposition and Shannon--Lupanov lookup on the resulting pieces give a circuit with at most \((1-\delta+o(1))t\) gates. To see the leading constant, pieces shorter than \(t/n^3\) have total length \(O(t/n^3)\), while all others have address arity \((1+o(1))\log_2t\). Prefix tests and multiplexing add \(O(n^2)\) gates. XOR with the given easy table then fits the hard threshold. Letting \(\delta\downarrow0\) gives both uniform assertions.


\section{Sharp threshold and rational degrees}
Use signs \(\{-1,1\}\) for the raw table entries. The threshold degree of the promise is the least degree of a polynomial having opposite strict signs on the two promised sets. For rational approximation we require a quotient with positive denominator and values in \([0,1]\) on the entire Boolean cube, with error at most \(1/3\) on the promise. Its degree is the maximum of the two polynomial degrees. Allowing a nonvanishing denominator of arbitrary sign changes the sign-degree lower comparison by at most a factor of two, on multiplying by that denominator.

\begin{theorem}
For \(s=\Theta(N^a),\ t=\Theta(N^b)\) with fixed \(0<a<b<1\), the threshold degree and the globally bounded rational approximate degree of gap MCSP are both \(\Theta_{a,b}(s)\).
\end{theorem}
\begin{proof}
Write \(M=|A_s|\). Apply the preceding minimax argument to \(A_s\cup(-A_s)\subseteq A_{s+1}\), accounting for a possibly charged output negation. This gives, for each hard table \(h\), a distribution \(\mu_h\) on positions with
\[
  \left|\sum_i\mu_h(i)h_i y_i\right|\le\epsilon
       \quad(y\in A_s),\qquad \epsilon=O(\sqrt{(s+1)n/t}).
\]
A disagreement lower bound also bounds agreement from below by comparison with the complemented circuit. Choose the constant in \(\epsilon\) large enough; for large \(n\) it is less than one. Take
\[
 d=\left\lceil\frac{\log(20M)}{2\log(1/\epsilon)}\right\rceil=O_{a,b}(s).
\]
The bound uses \(\log M=O(s\log s)\) and \(\log(1/\epsilon)=\Omega_{a,b}(\log N)\). Let \(\phi(x)\) be the vector of all squarefree parity monomials through degree \(d\), including the constant, and put
\[
  K=M^{-1}\sum_{y\in A_s}\phi(y)\phi(y)^{\mathsf T},\quad
  G=K+\epsilon^{2d}I,\quad
  P(x)=\phi(x)^{\mathsf T}G^{-1}\phi(x).
\]
All matrices are finite and \(G\) is positive definite. The polynomial \(P\) is nonnegative and has degree at most \(2d\). Each easy table \(y\) has \(P(y)\le M\), since \(G\succeq\phi(y)\phi(y)^{\mathsf T}/M\); conjugate by \(G^{-1/2}\).

For hard \(h\), expand
\[
     q_h(x)=\left(\sum_i\mu_h(i)h_ix_i\right)^d
            =v_h^{\mathsf T}\phi(x)
\]
on the cube. Before reduction the sum of absolute multinomial coefficients is \(1\). Using \(x_i^2=1\) and collecting terms gives \(\|v_h\|_2\le\|v_h\|_1\le1\).
Therefore \(v_h^{\mathsf T}Gv_h\le 2\epsilon^{2d}\), since \(q_h^2\le\epsilon^{2d}\) on every easy table. But \(q_h(h)=1\). Matrix Cauchy--Schwarz yields
\(P(h)\ge 1/(2\epsilon^{2d})\ge10M\).
Thus \(P-3M\) strictly separates the two classes, and \(3M/(P+3M)\) is a globally bounded rational approximant with the required orientation and accuracy.

For the lower bound use the one-mark construction above, with \(k=\Theta(t)\) buckets of size \(H=N/k\). The tuples of easy labels form a binary linear subgroup \(C\subseteq(\mathbb F_2^{\log_2H})^k\), since polynomial evaluation and coordinate projection are linear in their coefficients. Its logarithmic cardinality is \(O(sn)\) bits. Translate it by a uniformly random label tuple \(v\). Each fixed element of \(v+C\) is then a uniform label tuple, and a uniform tuple's marked table has size at most \(t\) with probability at most \(|A_t|/H^k\). The constant in \(k/t\) was chosen with fixed slack, so \(|C||A_t|/H^k=o(1)\). Some coset therefore consists entirely of hard tables. Both uniform laws on \(C\) and on this coset give independent uniform labels in every set of at most \(D=\Theta(s)\) buckets. Every raw-table monomial of degree at most \(D\) has identical expectations on the two laws, because it touches at most \(D\) buckets. A strictly sign-separating polynomial of that degree would have opposite expected signs. Finally, subtracting \(1/2\) from a rational approximant and clearing its positive denominator gives a sign polynomial. This proves both lower bounds.
\end{proof}
The positive quadratic form \(P\) can take large values away from the easy set. The rational denominator supplies the global boundedness used here; the proof gives a distinct upper bound for a different approximation model than bounded ordinary polynomials.

The bounded-error quantum query complexity with postselection is consequently
$\Theta(s)$. For the upper bound apply the amplitude construction of
Mahadev and de Wolf to the numerator and denominator just obtained.
Its success analysis is pointwise, so the approximation guarantee is needed
only on the promise, and the denominator is positive everywhere.
Conversely, if a $q$-query algorithm postselects with positive probability
on its promise, its joint acceptance and postselection probability $A(x)$
and its postselection probability $B(x)$ are polynomials of degree at most
$2q$. The polynomial $2A-B$ strictly separates the two promise classes,
which gives the lower bound even when postselection can fail identically
off promise. See U.~Mahadev and R.~de Wolf,
\emph{Rational approximations and quantum algorithms with postselection},
Quantum Information and Computation 15 (2015), 295--307.

\section{Optimal local bias and unrestricted threshold weight}
Let $N=2^n$. For a law $\nu$ on $N$-bit words write
$\beta(\nu)=\max_{\emptyset\ne S\subseteq[N]}|\mathbb E_\nu(-1)^{\sum_{i\in S}x_i}|$.
Fixing the random seed of a local generator must give a Boolean address circuit of size at most $s$. Classical small-bias generators give short seeds; Naor and Naor, \emph{Small-bias probability spaces}, SIAM J. Comput. 22 (1993), 838--856, and the local-generator approach of Cheraghchi, Kabanets, Lu and Myrisiotis, ICALP 2019, give relevant context. We count gates of a single output word as a function of its address and determine the leading exponential rate.
\begin{theorem}
If $\log s/\log n\to\infty$ and $s\log s\le(1-\eta)N$ for a fixed $\eta>0$, then
\[
  -\log_2\inf_{\operatorname{supp}\nu\subseteq A_s}\beta(\nu)
     =(1/2+o(1))s\log_2 s .
\]
The construction also gives the lower bound on this negative logarithm whenever $s\log s=O(N)$ in the same growth range. All gate estimates use the counted De Morgan basis.
\end{theorem}
\begin{proof}
Fix a sufficiently small constant $\delta>0$. Put $k=\lfloor(1-\delta)\log_2s\rfloor$, $K=2^k$, $m=\lfloor(1-\delta)k\rfloor$ and $q=2^m$. Use a strong seeded extractor on $n$ bits for min-entropy $k$, with error $\delta$, output $m$ bits and $D=\operatorname{poly}_\delta(n)$ seeds, computable in $\operatorname{poly}_\delta(n)$ gates (including any hardwired preprocessing). These parameters follow, for example, from D. Doron and J. Ribeiro, \emph{Nearly-linear time seeded extractors with short seeds}, ECCC TR24-176, revision 2 (2025). Write $h_j$ for the map at seed $j$. For any set $S$ of at least $K$ addresses, the pairs $(j,h_j(i))$ over $i\in S$ occupy at least $(1-\delta)Dq$ distinct bins: the joint seed-output law on a uniform element of $S$ is within $\delta$ of uniform.

Set $\ell=n^2$ and use $R=\lfloor (1-\delta)s m/(2\ell Dq)\rfloor$ independent copies of this array of bins. Thus the total number of bins $B=RDq$ satisfies $B=(1-\delta+o(1))sm/(2\ell)$, since $s/q$ dominates every fixed polynomial in $n$. In each bin $(a,j,u)$ choose two independent uniform elements $w_{a,j,u},v_{a,j,u}$ of $\mathbb F_{2^\ell}$. The output bit at the address interpreted as the integer $0\le i<N$ is the sum over copies $a$ and extractor seeds $j$ of
\[
  \operatorname{Tr}(w_{a,j,h_j(i)}v_{a,j,h_j(i)}^i),
\]
plus an independent $K$-wise uniform bit at $i$. For the latter evaluate one random polynomial of degree $K-1$ on the distinct addresses in $\mathbb F_{2^n}$ and project to one nonzero binary linear functional. This costs $O(K\operatorname{poly}(n))=o(s)$ gates by Horner evaluation. Use $v^0=1$ also for $v=0$.

For fixed $a,j$, the $2\ell$ seed bit tables indexed by $u$ each have $q$ bits and, by ordinary single-output Lupanov synthesis, cost at most $(1+o(1))q/m$ gates each. The resulting total is $(1+o(1))2\ell B/m$. All $RD=B/q$ trace evaluations by repeated squaring, all extractor evaluations, and the final XOR cost $O_\delta((B/q+D)\operatorname{poly}_\delta(n))=o(s)$. This verifies the size cap even with counted negations.

Every nonempty $S$ of size at most $K$ has zero bias by the independent summand. For larger $S$, every occupied bin has a nonempty set $S'$ of distinct addresses. Averaging its trace sign over $w,v$ gives the probability that $\sum_{i\in S'}v^i=0$. This is at most $2^{n-\ell}$ by the univariate root bound. Independence across the bins gives total bias at most $2^{-(\ell-n)(1-\delta)B}$. Sending $\delta$ to zero after the asymptotic limit proves the construction bound. Conversely Fourier Parseval applied to the density $2^N\nu$ yields
\[
 \beta(\nu)^2\ge\frac{2^N/|A_s|-1}{2^N-1}.
\]
The rooted census $\log_2|A_s|\le(1+o(1))s\log_2s$ finishes the stated equality.
\end{proof}

For a threshold polynomial use signs and require margin at least one on the promise. Its weight is the sum of absolute real coefficients in the multilinear Walsh basis; write $W_d$ for the minimum weight at degree at most $d$. The unrestricted degree is $N$.
\begin{theorem}
Under the preceding growth assumptions, if $s\le t$ and
$|A_t|2^{-N}=o(2^{-s\log_2 s/2})$, the unrestricted minimum weight for distinguishing $A_s$ from the complement of $A_t$ satisfies
\[
 \log_2 W_N=(1/2+o(1))s\log_2s.
\]
For polynomially separated $s=N^{a+o(1)}$, $t=N^{b+o(1)}$, $0<a<b<1$, one also has $\log_2 W_{Cs}=\Theta(s\log s)$ for sufficiently large constant $C$ whenever the bounded-ratio gate assumptions of the rational-degree theorem hold.
\end{theorem}
\begin{proof}
The function $2\mathbf1_{A_s}-1$ separates the promise. Parseval and Cauchy--Schwarz show that the Walsh weight of $\mathbf1_{A_s}$ is at most $\sqrt{|A_s|}$, giving the upper bound. For the converse compare the constructed generator with the uniform law conditioned on $A_t^c$. Every nonempty parity has conditional bias at most $|A_t|/(2^N-|A_t|)$. The difference of expectations of a separating polynomial is at least two, and is at most its weight times the sum of these two bias bounds. The constant coefficient cancels. For the low-degree upper in the polynomial-gap case use the preceding Gram separator divided by $M=|A_s|$; its feature dimension and inverse ridge are both at most $2^{O(s\log s)}$, so its Walsh weight is also at most that size after rescaling to unit margin.
\end{proof}

For comparison, the one-hot code gives a degree-dependent lower bound without a seeded extractor. If it matches independent uniform labels on $D$ buckets of size $H$, its hard coset has the same property. For a character testing at most $d$ raw cells let $Y$ count tested marks. The two laws match the binomial moments of $Y$ through $D$ and $\mathbb E\binom YD\le\binom dD H^{-D}$. Indeed $D$ designated cells can all be marked only when in distinct buckets. Truncate $(-1)^Y=\sum_j(-2)^j\binom Yj$ just before $j=D$. The error vanishes for $Y<D$ and is at most $(1+3^D)\binom YD$ otherwise, since $\binom Yj\le\binom YD\binom Dj$ for $j\le D\le Y$. Thus every margin-one threshold polynomial of degree $d\ge D$ has weight at least $(DH/(Cd))^D$ for an absolute $C$.

\section{Computing MCSP from XOR shares}
Alice receives a full word $X$, Bob receives $Y$, and the promised input to MCSP is $X\oplus Y$. Write $T=s\log_2 s$. Suppose $\log s/\log n\to\infty$, $T=o(N)$, $s\le t$, and $|A_t|=o(2^N)$. Then bounded-error interactive randomized communication is $(1+o(1))T$ bits. Interactive quantum communication with prior entanglement is $(1/2+o(1))T$ qubits. One-way quantum communication without prior entanglement, even with public randomness, is $(1+o(1))T$ qubits.

For the upper bounds Alice sends $\lceil\log_2|A_s|\rceil+2$ public-coin parity fingerprints. Bob accepts if some easy word agrees with the syndrome of $X\oplus Y$. The union bound gives error at most $1/4$, including for exact MCSP when $t=s$. Entanglement-assisted superdense coding sends these bits in half as many qubits to within rounding.

For the lower bounds compare the laws $(X,X\oplus Z)$ for uniform $X$ independent of $Z$, with $Z$ from the local generator or uniform. The latter violates the hard promise with probability $o(1)$. The bias of the generator is $\beta=2^{-(1/2-o(1))T}$. Let $U$ denote uniform measure on words, and write $\widehat f(S)=\mathbb E_U f\chi_S$. For vector-valued square-integrable maps $u,v$, independence and Parseval bound the change in expectation of $\langle u(X),v(Y)\rangle$ between these two worlds by
\[
 \beta\Big(\mathbb E_U\|u\|_2^2\Big)^{1/2}
       \Big(\mathbb E_U\|v\|_2^2\Big)^{1/2}.
\]
Apply this first to the indicators of the sides of a leaf rectangle of a deterministic $c$-bit protocol. The bound is $\beta\sqrt{U(A)U(B)}$; summing over the at most $2^c$ disjoint rectangles gives at most $\beta2^{c/2}$ by Cauchy--Schwarz. Public coins are mixtures and obey the same bound. For an arbitrary interactive $c$-qubit protocol, Linial and Shraibman's factorization theorem for its acceptance matrix gives vectors with the product of their maximum norms at most $2^c$, with or without prior entanglement (see A. Sherstov, \emph{Strong direct product theorems for quantum communication and query complexity}, Theorem 2.6). The preceding Fourier bound is then $\beta2^c$.

For one-way communication without entanglement take the two vector factors to be the density matrix $\rho_x$ of the $c$-qubit message and Bob's acceptance effect $0\le M_y\le I$, using the Hilbert--Schmidt inner product. Their norms are at most $1$ and $2^{c/2}$ respectively. This gives the asserted sharper one-way lower bound. In every case bounded-error correctness requires a constant distinction.

\paragraph{Unbounded error.}
Suppose additionally $|A_t|/2^N=o(2^{-T/2})$. The sign-rank of every completion of the XOR-share promise is at least $2^{(1/2-o(1))T}$, and one completion has sign-rank at most $2|A_s|+1$. Thus private-coin unbounded-error communication is between $(1/2-o(1))T$ and $(1+o(1))T$ bits by the sign-rank characterization of R. Paturi and J. Simon, \emph{Probabilistic communication complexity}, J. Comput. System Sci. 33 (1986), 106--123. See also A. Chattopadhyay and N. Mande, \emph{A short list of equalities induces large sign-rank}, SIAM J. Comput. 51 (2022), on a different nonsymmetric XOR function.

Indeed, start with $J_{xy}=2^N\nu(x\oplus y)-1$. Its operator norm is at most $2^N\beta$, by Fourier diagonalization and cancellation of the constant term. On easy pairs raise the entry to $1$ when needed. On unspecified pairs choose the prescribed completion sign with magnitude one. Each correction has magnitude at most $2$ since the law is supported on $A_s$. Every row and every column contains at most $|A_t|$ changed entries, whether or not the completion is translation invariant. The operator norm consequently increases by at most $2|A_t|$. The generalized Forster bound gives sign-rank at least $2^N/\|J_{\mathrm{corrected}}\|_{\mathrm{op}}$ since the minimum magnitude is one or greater. For the upper bound fix any injection $I$ from words to integers. The sign of
\[
 \frac12-\prod_{e\in A_s}(I(x)-I(y\oplus e))^2
\]
is positive precisely on the easy pairs, and the polynomial in $I(x)$ has at most $2|A_s|+1$ powers.

For the full exact matrix the exponent $T/2$ is optimal for this whole-matrix spectral certificate. In fact any sign-compatible witness of minimum magnitude one can be averaged over simultaneous translations without increasing operator norm or reducing the minimum magnitude. Its resulting entries have the form $g(x\oplus y)$. If its normalized norm is $\epsilon$, then every Walsh coefficient of $g$ is at most $\epsilon$ in magnitude. Writing $p=|A_s|/2^N$, the negative entries on $A_s^c$ give
$1-p-\epsilon\le\mathbb E_U g\mathbf1_{A_s}\le\epsilon\|\widehat{\mathbf1}_{A_s}\|_1\le\epsilon\sqrt{|A_s|}$. Improving sign-rank by a leading factor in the exponent would thus require further information.


\section{Scope}
Oliveira, Pich and Santhanam study circuit-size gaps in their hardness-magnification theorems \cite{ops}. Their hypotheses concern the size of computational models for deciding the promise problem. The results here concern oracle access to the input. In particular, the finite-family upper bounds may enumerate exponentially many circuits between queries. Our proof for a growing size ratio gives lower bounds on randomized and quantum algorithms. Establishing an equally sharp ordinary bounded approximate-degree lower bound in that growing-gap regime remains open in this work.
\begin{thebibliography}{9}
\bibitem{rank} S. Ballet and J. Pieltant. \href{https://arxiv.org/abs/1003.1864}{On the tensor rank of multiplication in any extension of $\mathbb F_2$}. Finite Fields and Their Applications, 2011.
\bibitem{kors} A. D. Korshunov. \href{https://www.mathnet.ru/eng/rm9459}{Computational complexity of Boolean functions}. Russian Math. Surveys 67 (2012), Section 2.4.
\bibitem{ops} I. C. Oliveira, J. Pich and R. Santhanam. \href{https://theoryofcomputing.org/articles/v017a011/}{Hardness magnification near state-of-the-art lower bounds}. Theory of Computing 17 (2021), article 11.
\end{thebibliography}
\end{document}
